\chapter{Estado del arte}

Se han revisado diversos trabajos relacionados con el reconocimiento y la traducción de lenguas de señas, incluyendo sistemas que utilizan visión artificial, aprendizaje profundo y procesamiento de lenguaje natural.

\section {Gramática de la lengua de señas mexicana \cite{ref42}}

La tesis de Miroslava Cruz Aldrete (2008) realiza una descripción detallada de la gramática de la Lengua de Señas Mexicana (LSM), abordando desde aspectos fonológicos hasta estructuras sintácticas complejas. Este estudio representa un esfuerzo significativo por documentar y analizar una lengua visogestual que ha sido históricamente poco explorada, ofreciendo datos sólidos que pueden servir como base para investigaciones futuras en el campo de la lingüística aplicada a las lenguas de señas. Entre los aportes más destacados, se encuentra la exposición de las características fonológicas y morfosintácticas de la LSM, así como la relevancia de los elementos no manuales en la construcción de significado, reafirmando que estas lenguas poseen una estructura gramatical compleja y autónoma comparable a la de cualquier lengua oral.

\section{Dispositivo traductor de LSM a español escrito para la extremidad superior derecha \cite{ref8}}

Este proyecto se enfoca en el diseño e implementación de un prototipo de dispositivo que se coloca en la extremidad superior derecha para traducir la Lengua de Señas Mexicana (LSM) a español escrito. Este dispositivo es capaz de traducir tanto el alfabeto dactilológico como un conjunto de señas básicas de la LSM. Para lograr esto, emplea una red neuronal artificial que clasifica las señas. Además, utiliza una combinación de sensores para identificar la posición y el movimiento de la mano y el brazo, incluyendo sensores resistivos para medir la flexión y un acelerómetro y giroscopio para medir la orientación y el movimiento.  Este proyecto fue desarrollado por Casa Sánchez, M. del R., en la UPIITA IPN.

\section{Reconocimiento de las señas estáticas del LSM con características basadas en aprendizaje profundo \cite{ref23}}

Este proyecto se enfoca en el reconocimiento de 21 señas estáticas de la Lengua de Señas Mexicana (LSM) utilizando características extraídas por la arquitectura VGG16 y clasificadores Random Forest (RF) y Support Vector Machines (SVM). Se evalúan las características profundas extraídas en las capas fc6 y fc7 de VGG16, así como la concatenación de ambas, para identificar las señas estáticas del LSM. El proyecto destaca por el uso de MediaPipe para la extracción de características de las manos, lo que permite una captura precisa y eficiente de los datos. Además, analiza las características extraídas por VGG16 en distintas capas y su combinación para determinar las más discriminativas para el reconocimiento de señas. Finalmente, utiliza clasificadores RF y SVM para el reconocimiento de señas, demostrando su eficacia en esta tarea. Este proyecto fue desarrollado por Rafael Fernández Rodríguez, Francisco Javier Peralta Rosas, Luis Ángel Zuñiga-Madrid, y Pedro Arguijo en el Tecnológico Nacional de México, Campus Misantla.

\section{Sistema de interpretación de imágenes para el reconocimiento y traducción del lenguaje de señas \cite{ref24}}

Este trabajo presenta el desarrollo de un sistema de reconocimiento de gestos que traduce imágenes a representaciones en lenguaje de señas. Utiliza OpenCV y Python para procesar imágenes, detectar y segmentar manos, y reconocer gestos predefinidos. Aunque no se enfoca específicamente en la LSM, demuestra el uso de técnicas de visión por computadora para la interpretación de señas. Emplea OpenCV para el procesamiento de imágenes y la detección de manos, lo que proporciona una base para la captura y análisis de gestos. Además, implementa técnicas de segmentación para aislar las manos del fondo de la imagen, mejorando la precisión del reconocimiento de gestos. Finalmente, utiliza algoritmos de reconocimiento de patrones para identificar gestos específicos a partir de las características extraídas de las manos.  Este proyecto fue desarrollado por Dora María Calderón Nepamuceno, Gabriela Kramer Bustos, y Efrén González Gómez en la Universidad Autónoma del Estado de México, Centro Universitario Nezahualcóyotl.

\section{Desarrollo de un intérprete de 20 frases en lengua de señas ecuatoriana a voz en tiempo real usando redes neuronales recurrentes \cite{ref25}}

Este proyecto se centra en el desarrollo de un intérprete en tiempo real que traduce 20 frases de la Lengua de Señas Ecuatoriana (LSEC) a voz. Su objetivo principal es facilitar la comunicación entre personas sordas y oyentes, mejorando la accesibilidad e inclusión de la comunidad sorda. El sistema se caracteriza por procesar y traducir las señas a voz de manera instantánea, enfocarse específicamente en la Lengua de Señas Ecuatoriana, y emplear Redes Neuronales Recurrentes (RNNs), en particular una combinación de CNN y LSTM, para analizar y comprender las secuencias de gestos que componen las frases en LSEC.  Desarrollado por Guevara Sanandrés, Juan Diego en la EPN.

\section{Las redes neuronales convolucionales en el uso de lenguaje de señas en la gestión académica de los estudiantes para SENATI \cite{ref26}}

Esta tesis explora la aplicación de redes neuronales convolucionales (CNN) para el reconocimiento del lenguaje de señas en el contexto de la gestión académica de estudiantes en el Servicio Nacional de Adiestramiento en Trabajo Industrial (SENATI). Se destacan dos aspectos principales: el empleo de CNN, un tipo de red neuronal profunda especializada en el procesamiento de imágenes, para analizar y reconocer los gestos del lenguaje de señas; y el enfoque en la aplicación práctica de la tecnología en procesos académicos como la atención al estudiante y consultas sobre especialidades. Desarrollado por Hugo Alejandro Mamanchura Lima en la UPLA.

\section{Diseño e implementación de un sistema de traducción de lenguaje de señas dactilológico a lenguaje natural escrito para personas no hablantes mediante visión artificial \cite{ref27}}

Este proyecto se centra en el diseño e implementación de un sistema que utiliza visión artificial para traducir el lenguaje de señas dactilológico (deletreo con los dedos) al español escrito, con el objetivo de mejorar la comunicación de las personas no hablantes.  El sistema emplea técnicas de visión artificial para capturar y procesar los gestos de las manos en tiempo real.  Utiliza la biblioteca MediaPipe para la detección y seguimiento de manos, lo que permite identificar 21 puntos de referencia en la mano para un análisis preciso. Además, implementa una Red Neuronal Convolucional (CNN) para clasificar los gestos de las manos y traducirlos a letras del alfabeto. Desarrollado por Cusme Quintanilla, Medelin Lisset en la Escuela Superior Politécnica de Chimborazo.

\section{Reconocimiento de señas de la Lengua de Señas Mexicana mediante técnicas de Machine Learning \cite{ref28}}

Este artículo evalúa el rendimiento de diferentes técnicas de Machine Learning (Aprendizaje Automático) para reconocer señas del alfabeto dactilológico de la Lengua de Señas Mexicana (LSM).  Se recolectó un conjunto de datos de imágenes de señas y se entrenaron modelos utilizando cuatro técnicas.  El estudio se enfoca en el reconocimiento de señas del alfabeto dactilológico de la LSM, que no involucran movimiento.  Se comparan cuatro técnicas para el reconocimiento de señas: CNN, Transfer Learning, Teachable Machine y CNN con características específicas. Para el entrenamiento y la evaluación, se utilizó un conjunto de datos propio de imágenes de señas.  En una de las técnicas, se extrajeron características específicas de las imágenes (coordenadas de puntos de referencia en las manos) para mejorar el rendimiento del modelo. Desarrollado por Gabriel A. Salgado-Martínez, René E. Cuevas-Valencia, Angelino Feliciano-Morales, y Arnulfo Catalán-Villegas en la Escuela Superior de Tlahuelil.

\section{Desarrollo de un Sistema Web intérprete de la Lengua de Señas Argentina \cite{29}}

Esta tesis presenta el desarrollo de un sistema web que interpreta un conjunto de señas de la Lengua de Señas Argentina (LSA) en tiempo real o a partir de videos pregrabados. El sistema utiliza modelos de aprendizaje automático, incluyendo árboles de decisión optimizados (XGBoost) y redes neuronales recurrentes (RNN, específicamente LSTM y GRU), para traducir secuencias de imágenes (videos) a su correspondiente texto en español.  El sistema es capaz de procesar videos tanto en tiempo real como pregrabados, extrayendo información relevante de las secuencias de imágenes.  Para la detección y seguimiento de manos, utiliza la biblioteca MediaPipe, lo que facilita la extracción de características de los gestos. Desarrollado por Facundo Nicolas Recabarren Martin en la Universidad Nacional de San Juan, Facultad de Ciencias Exactas, Físicas y Naturales, Departamento de Informática.

\section{Desarrollo de algoritmos para Traductor automático de lenguaje de señas Mexicanas (LSM) \cite{ref30}}

Esta tesis doctoral se enfoca en el desarrollo de algoritmos para la traducción automática del Lenguaje de Señas Mexicano (LSM). Aborda el desafío de reconocer y clasificar señas a partir de secuencias de imágenes, considerando aspectos como el movimiento de las manos, la posición corporal y los puntos de contacto.  El análisis de videos de señas se realiza extrayendo 15 cuadros por palabra y enfocándose en regiones de interés, como las manos y la cara.  Se utilizan características geométricas, como área, redondez, longitud del borde, elongación, centro de gravedad y densidad, para describir las señas. Para la clasificación, se emplean y comparan diferentes algoritmos, incluyendo Support Vector Machines (SVM), árboles de decisión, clasificador bayesiano y redes neuronales, para evaluar su eficacia en el reconocimiento de señas. Desarrollado por Josué Espejel Cabrera en la Universidad Autónoma del Estado de México.


\section{Sistema multi-agente de reconocimiento de frases en LSM utilizando CBR \cite{ref31}}

Esta tesis presenta el desarrollo de un sistema multi-agente que reconoce frases en Lengua de Señas Mexicana (LSM) y las traduce a través de un avatar en un dispositivo móvil. El sistema utiliza la técnica de Razonamiento Basado en Casos (CBR) para mejorar su capacidad de comunicación al aprender de nuevas frases. Se enfoca en el reconocimiento de frases completas en LSM, en lugar de palabras aisladas, lo que permite una comunicación más natural y fluida.  Emplea una arquitectura multi-agente, donde cada agente tiene una tarea específica, como recibir la frase, procesarla y proyectar la traducción.  El uso de CBR permite almacenar y reutilizar experiencias pasadas, lo que permite al sistema aprender de nuevas frases y mejorar su precisión con el tiempo. Desarrollado por Ing. César René Martínez Aguirre en el Instituto Tecnológico de Hermosillo.


\section{Modelos Computacionales Aplicados a la Lengua de Señas Mexicana \cite{ref32}}

Esta tesis se enfoca en el desarrollo y evaluación de un sistema para reconocer la dactilología (alfabeto manual) y los primeros diez números de la Lengua de Señas Mexicana (LSM) utilizando visión computacional y técnicas de aprendizaje automático y profundo. Se creó una base de datos propia de LSM y se aplicó MediaPipe para la extracción de características de las manos, obteniendo 21 puntos de referencia por mano. Se entrenaron y compararon modelos como Support Vector Machine (SVM) para reconocer las señas estáticas y Gated Recurrent Unit (GRU) para el reconocimiento de señas continuas. Desarrollado por Ing. Mario Emmanuel Rodríguez Trejo en la Universidad Autónoma del Estado de Morelos.

\section{Modelo para la creación y uso de un repositorio digital de segmentos visuales en lengua de señas mexicanas utilizando Inteligencias Artificial \cite{ref33}}

Esta tesis propone un modelo para crear y utilizar un repositorio digital de segmentos visuales en Lengua de Señas Mexicana (LSM) usando inteligencia artificial. El objetivo es mejorar la plataforma de traducción español-LSM desarrollada previamente en el Instituto Tecnológico de Hermosillo, facilitando la comunicación entre personas sordas y oyentes en contextos educativos.  El repositorio almacena segmentos visuales de señas en LSM extraídos de videos disponibles en la web.  Utiliza técnicas de visión e inteligencia artificial para extraer automáticamente las señas y su  texto correspondiente de los videos.  Además, convierte los segmentos de video en representaciones 3D animadas utilizando el modelo SMPL, permitiendo una visualización más realista y versátil.  Desarrollado por Maria Luisa Galaz Palma en el Instituto Tecnológico de Hermosillo.

\section{Análisis Comparativo de Proyectos}

Para contextualizar el alcance del presente trabajo y contrastarlo con los estudios previos, se ha elaborado la \textbf{Tabla \ref{tab:comparativa_proyectos}} que resume las metodologías y características tecnológicas empleadas en los proyectos revisados. Este análisis comparativo permite identificar tendencias, debilidades comunes y el nivel de complejidad de las soluciones implementadas en el reconocimiento de lenguaje de señas, sirviendo como una base sólida para la justificación metodológica.

\begin{table}[htbp]
    \centering
    \caption{Comparativa de proyectos relacionados con el reconocimiento de lenguaje de señas}
    % Se usa \small para un mejor ajuste y legibilidad. Se corrige la doble línea.
    \small
    % Definición de la estructura de la tabla (solo una línea entre LSM y COB)
    \begin{tabular}{|p{4.5cm}|c|c|c|c|c|c|c|c|c|}
        \hline
        \textbf{Proyecto Revisado} &
        \textbf{CV} & % Captura Video/Img
        \textbf{TR} & % Tiempo Real
        \textbf{MP} & % MediaPipe
        \textbf{NUBE} & % Nube (AWS, etc.)
        \textbf{S(T/V)} & % Salida (Texto/Voz)
        \textbf{RN} & % Redes Neuronales
        \textbf{SEG} & % Segmentación/ROI
        \textbf{LSM} & % Enfoque en LSM
        \textbf{COB (\%)} \\ \hline
        % APLICO LA CORRECCIÓN DE LA SINTAXIS: \cellcolor[rgb]{R, G, B}
        \textbf{Sistema de reconocimiento de señas de detección de palabras para el lenguaje de señas LSM} & Sí & Sí & Sí & Sí & Sí & Sí & Sí & Sí & \textbf{\cellcolor[rgb]{0.8, 1.0, 0.8} 100\%} \\ \hline
        Dispositivo traductor de LSM a español         & Sí & No & No & No & Sí & Sí & No & Sí & \cellcolor[rgb]{1.0, 0.8, 0.6} 50\% \\ \hline
        Reconocimiento de señas estáticas del LSM      & Sí & No & Sí & No & No & Sí & No & Sí & \cellcolor[rgb]{1.0, 0.8, 0.6} 50\% \\ \hline
        Sistema de interpretación de imágenes          & Sí & Sí & No & No & No & No & Sí & No & \cellcolor[rgb]{1.0, 0.7, 0.7} 37.5\% \\ \hline
        Intérprete en tiempo real de LSEC              & No & Sí & No & No & Sí & Sí & No & No & \cellcolor[rgb]{1.0, 0.8, 0.6} 50\% \\ \hline
        Reconocimiento de señas en gestión académica   & No & No & No & No & No & Sí & No & No & \cellcolor[rgb]{1.0, 0.7, 0.7} 12.5\% \\ \hline
        Traducción de señas dactilológicas             & Sí & No & Sí & No & Sí & Sí & No & No & \cellcolor[rgb]{1.0, 0.8, 0.6} 50\% \\ \hline
        Reconocimiento de alfabeto dactilológico       & Sí & No & No & No & No & Sí & No & Sí & \cellcolor[rgb]{1.0, 0.7, 0.7} 25\% \\ \hline
        Sistema web intérprete de LSA                  & Sí & Sí & Sí & No & Sí & Sí & No & No & \cellcolor[rgb]{1.0, 1.0, 0.8} 62.5\% \\ \hline
        Traductor automático de LSM                    & Sí & No & No & No & No & Sí & Sí & Sí & \cellcolor[rgb]{1.0, 0.8, 0.6} 50\% \\ \hline
        Sistema multi-agente para frases en LSM        & No & Sí & No & No & Sí & No & No & Sí & \cellcolor[rgb]{1.0, 0.7, 0.7} 37.5\% \\ \hline
        Modelos computacionales aplicados a LSM        & Sí & No & Sí & No & No & Sí & No & Sí & \cellcolor[rgb]{1.0, 0.8, 0.6} 50\% \\ \hline
        Repositorio digital para segmentos visuales    & Sí & No & No & No & No & Sí & No & Sí & \cellcolor[rgb]{1.0, 0.7, 0.7} 37.5\% \\ \hline
    \end{tabular}

    % --- Explicación de Acrónimos ---
    \begin{flushleft}
        \footnotesize
        \vspace{1ex}
        \textbf{Acrónimos:} \\
        \textbf{CV:} Captura de Video/Imágenes. \textbf{TR:} Tiempo Real. \textbf{MP:} Uso de MediaPipe. \textbf{NUBE:} Uso de servicios en la Nube (AWS, Azure, etc.). \textbf{S(T/V):} Salida en Texto o Voz. \textbf{RN:} Uso de Redes Neuronales. \textbf{SEG:} Segmentación o uso de Región de Interés (ROI). \textbf{LSM:} Enfoque específico en la Lengua de Señas Mexicana. \textbf{COB (\%):} Porcentaje de Cobertura de las características listadas.

        \vspace{1ex}
        \textbf{Leyenda de Cobertura:} \
        \textcolor[rgb]{0.0, 0.5, 0.0}{\textbf{Verde (100\%):}} Cobertura Completa. \
        \textcolor[rgb]{0.8, 0.8, 0.0}{\textbf{Amarillo (62.5\%):}} Cobertura Alta. \
        \textcolor[rgb]{1.0, 0.5, 0.0}{\textbf{Naranja (50\%):}} Cobertura Media. \
        \textcolor[rgb]{0.8, 0.0, 0.0}{\textbf{Rojo (<50\%):}} Cobertura Baja.
    \end{flushleft}

    \label{tab:comparativa_proyectos}
\end{table}

